%%%%%%%% ICML 2026 SUBMISSION FILE %%%%%%%%%%%%%%%%%

\documentclass{article}

% Recommended packages
\usepackage{microtype}
\usepackage{graphicx}
\usepackage{subcaption}
\usepackage{booktabs}
\usepackage{hyperref}

% Attempt to make hyperref and algorithmic work together better:
\newcommand{\theHalgorithm}{\arabic{algorithm}}

% Use the following line for the initial blind version submitted for review:
\usepackage{icml2026}

% For camera-ready, use:
% \usepackage[accepted]{icml2026}

\usepackage{amsmath}
\usepackage{amssymb}
\usepackage{algorithm}
\usepackage{algorithmic}

% Custom commands
\newcommand{\deltacache}{\textsc{DeltaCache}}

% Short running title
\icmltitlerunning{DeltaCache: Layer-Aware KV Cache Management}

\begin{document}

\twocolumn[
\icmltitle{DeltaCache: Layer-Aware KV Cache Management Guided by Transformer Semantics}

\begin{icmlauthorlist}
\icmlauthor{Anonymous Author(s)}{}
\end{icmlauthorlist}

\icmlkeywords{Large Language Models, KV Cache, Inference Optimization, Prefix Caching}

\vskip 0.3in
]

\printAffiliationsAndNotice{}

\begin{abstract}
Large Language Models (LLMs) face significant computational overhead during inference due to repeated Key-Value (KV) cache computations for shared prompt prefixes.
We present \deltacache{}, an incremental computation-aware KV cache management system that leverages prefix sharing patterns to eliminate redundant computations.
Our key innovations include: (1) a prefix tree-based cache architecture for efficient KV reuse, (2) attention-aware eviction that uses attention scores to prioritize important cache entries, (3) layer-aware caching that prioritizes late-layer (semantic) over early-layer (syntactic) KV cache based on transformer interpretability insights, and (4) hierarchical GPU/CPU memory management with predictive prefetching.
On targeted application benchmarks (RAG, few-shot learning, code completion), \deltacache{} achieves \textbf{5.13$\times$} average speedup on Mistral-7B (fp16), with \textbf{10.90$\times$} cached speedup on code completion scenarios.
In comprehensive baseline comparisons, we achieve up to \textbf{13.43$\times$} speedup on 2K-token contexts.
Speedup scales at \textbf{2.43$\times$ per 1K additional tokens}, with memory overhead of only 17--22\%.
We verify 100\% correctness at fp16 precision and identify optimal scenarios: \deltacache{} excels in fixed-prefix workloads while multi-turn conversations show limited gains.
\end{abstract}

\section{Introduction}
\label{sec:intro}

Large Language Models (LLMs) have become foundational components in modern AI systems, powering applications from conversational agents to code generation \citep{brown2020language,openai2023gpt4}.
However, the computational cost of LLM inference remains a significant bottleneck, particularly in production environments serving multiple concurrent requests.

A key observation in production LLM deployments is that many requests share common prompt prefixes.
Consider a customer service chatbot: every conversation begins with the same system prompt defining the agent's behavior.
Similarly, Retrieval-Augmented Generation (RAG) systems often prepend the same document context to multiple user queries \citep{lewis2020retrieval}.
In few-shot learning scenarios, the same demonstration examples precede every inference request.

Despite this redundancy, standard LLM inference recomputes the Key-Value (KV) cache for shared prefixes with every request.
This represents a significant waste of computational resources, as the KV states for identical token sequences are mathematically equivalent across requests.
The quadratic complexity of self-attention makes this particularly costly for long prefixes.

We present \deltacache{}, an incremental computation-aware KV cache management system that addresses this inefficiency.
Our approach is grounded in two key insights from transformer interpretability research \citep{rogers2020primer,elhage2021mathematical}:

\begin{enumerate}
    \item \textbf{Prefix independence}: For causal language models, KV states for prefix tokens are independent of future tokens, enabling safe reuse.
    \item \textbf{Layer heterogeneity}: Different transformer layers serve distinct functions---early layers capture local syntactic patterns while late layers encode semantic features \citep{geva2023dissecting}---suggesting differentiated caching priorities.
\end{enumerate}

Our contributions are:
\begin{itemize}
    \item A \textbf{prefix tree-based cache architecture} that enables efficient lookup and reuse of cached KV states for shared prompt prefixes with $O(n)$ complexity.
    \item An \textbf{incremental computation engine} that computes only the delta (new tokens) while reusing cached KV states for matched prefixes.
    \item \textbf{Attention-aware eviction}: A novel policy that uses attention scores to identify and retain important cache entries.
    \item \textbf{Layer-aware caching}: Prioritizes late-layer (semantic) KV cache over early-layer (local pattern) cache based on transformer layer characteristics.
    \item \textbf{Hierarchical memory management}: GPU$\rightarrow$CPU tiering with proactive offloading and predictive prefetching.
    \item \textbf{Comprehensive evaluation} demonstrating 5.13$\times$ average speedup on Mistral-7B with 100\% correctness at fp16 precision.
\end{itemize}

\section{Related Work}
\label{sec:related}

\paragraph{KV Cache Memory Optimization}
The KV cache is a critical component in transformer inference, storing computed attention states to avoid redundant computation during autoregressive generation \citep{vaswani2017attention}.
Several works focus on reducing KV cache memory footprint through quantization \citep{liu2024kivi,hooper2024kvquant}, compression \citep{ge2023model}, or sparse attention patterns \citep{child2019generating,xiao2024efficient}.
H2O \citep{zhang2024h2o} identifies ``heavy hitter'' tokens that receive high attention and retains only these.
SnapKV \citep{snapkv2024} compresses KV cache by clustering similar attention patterns.
These approaches focus on memory reduction rather than computation reuse, which is complementary to our work.

\paragraph{Prefix Caching in Serving Systems}
vLLM \citep{kwon2023efficient} introduced automatic prefix caching as an optimization for shared prompt prefixes.
Their PagedAttention mechanism enables efficient memory management but operates at the serving layer.
SGLang \citep{zheng2023efficiently} further optimizes prefix sharing through RadixAttention with efficient tree-based prefix matching.
ChunkAttention \citep{ye2024chunkattention} introduces prefix-aware KV cache with two-phase partition.
\deltacache{} differs by providing a general-purpose library that can be integrated with any transformer implementation, offering fine-grained control over eviction policies and layer-aware prioritization.

\paragraph{Layer-wise KV Cache Management}
Recent work recognizes that different transformer layers have different importance for KV caching.
LayerCondensed \citep{layercondensed2024} shares KV cache across adjacent layers.
SqueezeAttention \citep{squeezeattention2025} allocates layer-wise optimal budget for KV cache.
Our layer-aware caching is motivated by similar insights but focuses on eviction priority rather than compression.

\section{Methodology}
\label{sec:method}

\subsection{Problem Formulation}

Consider a transformer model with $L$ layers processing a sequence of tokens $\mathbf{x} = [x_1, x_2, \ldots, x_n]$.
For each layer $l$ and position $i$, the model computes key and value states:
\begin{align}
    \mathbf{k}_i^l &= W_K^l \mathbf{h}_i^{l-1} \\
    \mathbf{v}_i^l &= W_V^l \mathbf{h}_i^{l-1}
\end{align}
where $\mathbf{h}_i^{l-1}$ is the hidden state from the previous layer and $W_K^l, W_V^l$ are learned projection matrices.

The key insight enabling prefix caching is that for causal language models, the KV states for prefix tokens are \textit{independent of future tokens}.
Formally, $\mathbf{k}_i^l$ and $\mathbf{v}_i^l$ depend only on $[x_1, \ldots, x_i]$, not on any tokens after position $i$.
This property follows from the causal attention mask that prevents positions from attending to future tokens.

\textbf{Implication}: If we have previously computed KV states for a prefix $[x_1, \ldots, x_m]$, and we receive a new sequence $[x_1, \ldots, x_m, x_{m+1}, \ldots, x_n]$ with the same prefix, we can reuse the cached KV states for positions $1$ through $m$ and only compute new states for positions $m+1$ through $n$.

\subsection{System Architecture}

\deltacache{} consists of three main components: a prefix tree for cache organization, a memory pool for hierarchical storage, and an incremental computation engine.

\paragraph{Prefix Tree}
We organize cached KV states in a prefix tree (trie) data structure.
Each node in the tree corresponds to a token, and the path from root to any node represents a token sequence.
Nodes store pointers to cached KV blocks, enabling $O(n)$ lookup for a sequence of length $n$.

The prefix tree supports three operations:
\begin{itemize}
    \item \texttt{lookup(tokens)}: Find the longest matching prefix and return cached KV states. Returns (prefix\_length, kv\_cache).
    \item \texttt{insert(tokens, kv\_cache)}: Add a new sequence with its computed KV states to the tree.
    \item \texttt{evict(node)}: Remove a node and free associated memory when capacity is exceeded.
\end{itemize}

\paragraph{Memory Pool}
The memory pool manages GPU and CPU memory for KV cache blocks.
It implements a two-tier caching hierarchy:
\begin{itemize}
    \item \textbf{GPU tier}: Fast access for active cache entries
    \item \textbf{CPU tier}: Overflow storage for less frequently used entries
\end{itemize}
Proactive offloading moves entries from GPU to CPU before GPU memory is exhausted (default threshold: 80\%).
Predictive prefetching moves hot entries back to GPU based on access patterns.

\paragraph{Incremental Computation Engine}
The incremental engine coordinates the computation flow as shown in Algorithm~\ref{alg:incremental}.

\begin{algorithm}[t]
\caption{Incremental KV Computation}
\label{alg:incremental}
\begin{algorithmic}[1]
\REQUIRE Input tokens $\mathbf{x} = [x_1, \ldots, x_n]$, Model $M$, Cache $C$
\ENSURE Complete KV cache for full sequence
\STATE $(\text{prefix\_len}, \text{kv\_cached}) \leftarrow C.\text{lookup}(\mathbf{x})$
\IF{prefix\_len $> 0$}
    \STATE $\mathbf{x}_\text{delta} \leftarrow [x_{\text{prefix\_len}+1}, \ldots, x_n]$
    \STATE $\text{kv\_delta} \leftarrow M.\text{forward}(\mathbf{x}_\text{delta}, \text{kv\_cached})$
    \STATE $\text{kv\_full} \leftarrow \text{concat}(\text{kv\_cached}, \text{kv\_delta})$
\ELSE
    \STATE $\text{kv\_full} \leftarrow M.\text{forward}(\mathbf{x})$
\ENDIF
\STATE $C.\text{insert}(\mathbf{x}, \text{kv\_full})$
\STATE \textbf{return} kv\_full
\end{algorithmic}
\end{algorithm}

\subsection{Eviction Policies}

When memory is exhausted, \deltacache{} must evict cached entries.
We implement several eviction policies, including novel attention-aware and layer-aware strategies.

\paragraph{Attention-Aware Eviction}
We introduce \emph{attention-aware eviction}, which leverages attention scores to prioritize important cache entries.
The intuition is that KV cache entries receiving high attention from subsequent queries are likely important for future queries.

We compute entry importance as:
\begin{equation}
    I(e) = w_a \cdot \bar{\alpha}_e + w_f \cdot \log(1 + f_e) + w_r \cdot \frac{1}{1 + \log(1 + t_e)}
\end{equation}
where $\bar{\alpha}_e$ is the exponential moving average (EMA) of attention scores received by entry $e$, $f_e$ is access frequency, $t_e$ is time since last access, and $w_a, w_f, w_r$ are tunable weights (default: 0.5, 0.3, 0.2).

\paragraph{Layer-Aware Caching}
Research in transformer interpretability reveals that different layers serve distinct functions \citep{rogers2020primer,elhage2021mathematical}:
\begin{itemize}
    \item \textbf{Early layers} (0--30\%): Capture local syntactic patterns, positional information
    \item \textbf{Middle layers} (30--70\%): Build compositional representations
    \item \textbf{Late layers} (70--100\%): Encode task-specific semantic features
\end{itemize}

We assign layer-dependent importance weights using a sigmoid function:
\begin{equation}
    w_l = \frac{1}{1 + e^{-k(l/L - \tau)}}
\end{equation}
where $l$ is the layer index, $L$ is total layers, $k=5$ controls steepness, and $\tau=0.3$ is the transition point.
This gives early layers weight $\approx$0.18 and late layers weight $\approx$0.88, prioritizing semantic representations during eviction.

\subsection{Correctness Guarantee}

A critical requirement for \deltacache{} is numerical correctness.
The cached KV states must be identical (within floating-point tolerance) to freshly computed states.

We ensure correctness through:
\begin{enumerate}
    \item \textbf{Position-aware computation}: When computing delta tokens, we provide correct absolute position IDs accounting for the cached prefix length.
    \item \textbf{RoPE handling}: For models using Rotary Position Embeddings \citep{su2024roformer}, position IDs must match absolute positions in the full sequence. We explicitly pass \texttt{position\_ids} to the model forward pass.
    \item \textbf{Validation}: We verify $\|\text{kv}_\text{cached} - \text{kv}_\text{fresh}\|_\infty < \epsilon$ with tolerance $\epsilon = 2 \times 10^{-2}$.
\end{enumerate}

\subsection{Theoretical Analysis}

\paragraph{Complexity Reduction}
For a transformer with $L$ layers processing $n$ tokens with hidden dimension $d$:
\begin{itemize}
    \item \textbf{Baseline prefill}: $O(Ln^2d)$ for self-attention
    \item \textbf{With cached prefix $m$}: $O(L(n-m)^2d + Lmd)$
\end{itemize}

The theoretical speedup is:
\begin{equation}
    \text{Speedup} = \frac{n^2}{(n-m)^2 + O(m)} \approx \frac{n^2}{(n-m)^2}
\end{equation}

For $n=1000$ tokens with cached prefix $m=900$:
$\text{Speedup} \approx 1000^2 / 100^2 = 100\times$.

This quadratic relationship explains why \deltacache{} benefits increase dramatically with prefix length.

\paragraph{Memory Overhead}
The prefix tree introduces $O(V \cdot b)$ additional memory where $V$ is number of cached sequences and $b$ is average branching factor. In practice, this is negligible ($<$1\%) compared to KV cache storage.

\section{Experiments}
\label{sec:experiments}

\subsection{Experimental Setup}

\paragraph{Models}
We evaluate on two models representing different scales:
\begin{itemize}
    \item \textbf{TinyLlama-1.1B} \citep{tinyllama}: A compact Llama-architecture model
    \item \textbf{Mistral-7B} \citep{jiang2023mistral}: A 7B parameter model at fp16 precision
\end{itemize}

\paragraph{Hardware}
2$\times$ NVIDIA RTX 3090 GPUs (24GB each), 64GB CPU memory.

\paragraph{Metrics}
We report: (1) \textbf{Correctness}: Max KV difference and match rate; (2) \textbf{Speedup}: Baseline/\deltacache{} latency ratio; (3) \textbf{Token reuse}: Fraction of tokens served from cache; (4) \textbf{Cached speedup}: Speedup excluding first (cold) query.

\subsection{Correctness Validation}

Table~\ref{tab:correctness} validates numerical correctness across precisions.

\begin{table}[t]
\centering
\caption{Correctness validation (tolerance $2\times10^{-2}$)}
\label{tab:correctness}
\small
\begin{tabular}{lccc}
\toprule
Model & Precision & Max Diff & Accuracy \\
\midrule
TinyLlama-1.1B & fp16 & 0.0156 & 100\% \\
Mistral-7B & fp16 & 0.0156 & 100\% \\
Mistral-7B & 8-bit & 0.60 & 10\% \\
\bottomrule
\end{tabular}
\end{table}

At fp16, both models achieve \textbf{100\% correctness} with max difference $\sim$0.016, within expected fp16 precision.
The 8-bit quantized model shows larger differences due to quantization noise independent of caching.

\subsection{Targeted Application Benchmarks}

We design benchmarks for workloads where prefix caching is most beneficial: fixed long prefix combined with varying short queries.
This pattern is common in production deployments.

\paragraph{Scenarios}
\begin{itemize}
    \item \textbf{RAG with Fixed KB}: Python docs, ML tutorials, API references ($\sim$650 tokens prefix)
    \item \textbf{Few-shot Sentiment}: 5 labeled examples + test inputs ($\sim$280 tokens)
    \item \textbf{Long System Prompt}: Detailed chatbot instructions ($\sim$480 tokens)
    \item \textbf{Code Completion}: Large code context + completion ($\sim$1450 tokens)
\end{itemize}

\begin{table}[t]
\centering
\caption{Targeted benchmarks on Mistral-7B (fp16)}
\label{tab:targeted_mistral}
\small
\begin{tabular}{lcccc}
\toprule
Scenario & Prefix & Speedup & Cached & Reuse \\
\midrule
Code Completion & 1448 & \textbf{6.62$\times$} & \textbf{10.90$\times$} & 92.3\% \\
RAG: API Docs & 646 & 5.90$\times$ & 6.39$\times$ & 92.4\% \\
RAG: ML Tutorial & 647 & 5.56$\times$ & 6.46$\times$ & 92.3\% \\
RAG: Python Guide & 676 & 5.51$\times$ & 3.11$\times$ & 92.4\% \\
Long System Prompt & 484 & 3.83$\times$ & 4.74$\times$ & 92.2\% \\
Few-shot Sentiment & 282 & 3.39$\times$ & 3.93$\times$ & 89.4\% \\
\midrule
\textbf{Average} & -- & \textbf{5.13$\times$} & \textbf{5.92$\times$} & 91.8\% \\
\bottomrule
\end{tabular}
\end{table}

\begin{table}[t]
\centering
\caption{Targeted benchmarks on TinyLlama-1.1B}
\label{tab:targeted_tinyllama}
\small
\begin{tabular}{lcccc}
\toprule
Scenario & Prefix & Speedup & Cached & Reuse \\
\midrule
Code Completion & 1464 & \textbf{3.64$\times$} & \textbf{4.22$\times$} & 93.8\% \\
RAG: API Docs & 660 & 3.94$\times$ & 2.28$\times$ & 94.4\% \\
RAG: ML Tutorial & 656 & 2.96$\times$ & 2.13$\times$ & 94.2\% \\
Long System Prompt & 497 & 1.80$\times$ & 1.83$\times$ & 94.0\% \\
Few-shot Sentiment & 298 & 1.22$\times$ & 1.20$\times$ & 90.8\% \\
\midrule
\textbf{Average} & -- & \textbf{2.60$\times$} & \textbf{2.09$\times$} & 93.4\% \\
\bottomrule
\end{tabular}
\end{table}

Tables~\ref{tab:targeted_mistral} and \ref{tab:targeted_tinyllama} show results.
\deltacache{} achieves \textbf{5.13$\times$} average speedup on Mistral-7B with \textbf{10.90$\times$} cached speedup on code completion.
Larger models benefit more from caching due to higher per-token computation cost.

\subsection{Prefix Length Scaling}

\begin{figure}[t]
\centering
\includegraphics[width=\columnwidth]{../experiments/figures/prefix_length_scaling.pdf}
\caption{Speedup scaling with prefix length on TinyLlama-1.1B. Cached speedup reaches 5$\times$ at 1000 tokens.}
\label{fig:prefix_scaling}
\end{figure}

Figure~\ref{fig:prefix_scaling} and Table~\ref{tab:prefix_ablation} show how speedup scales with prefix length.

\begin{table}[t]
\centering
\caption{Prefix length ablation (TinyLlama-1.1B)}
\label{tab:prefix_ablation}
\small
\begin{tabular}{cccl}
\toprule
Prefix & Overall & Cached & Finding \\
\midrule
100 & 0.94$\times$ & 1.14$\times$ & Overhead $>$ benefit \\
250 & 1.98$\times$ & 2.14$\times$ & Breakeven point \\
500 & 2.77$\times$ & 3.21$\times$ & Clear benefit \\
750 & 3.12$\times$ & 3.76$\times$ & Strong benefit \\
1000 & \textbf{3.96$\times$} & \textbf{5.05$\times$} & Maximum tested \\
\bottomrule
\end{tabular}
\end{table}

Key finding: \textbf{Prefix length $>$200 tokens is required for benefit}.
At 100 tokens, cache management overhead exceeds computation savings.
Speedup grows with prefix length following the theoretical quadratic relationship.

\subsection{Scenario Comparison}

\begin{figure}[t]
\centering
\includegraphics[width=\columnwidth]{../experiments/figures/scenario_comparison.pdf}
\caption{Speedup comparison across targeted scenarios. Mistral-7B consistently outperforms TinyLlama due to higher baseline computation cost.}
\label{fig:scenario}
\end{figure}

Figure~\ref{fig:scenario} compares speedups across scenarios for both models.
The results reveal two key patterns:
\begin{enumerate}
    \item \textbf{Model size matters}: Larger models benefit more from caching (5.13$\times$ vs 2.60$\times$ average)
    \item \textbf{Prefix length matters}: Code completion (1448 tokens) achieves highest speedup (6.62$\times$)
\end{enumerate}

\subsection{Speedup vs Prefix Length}

\begin{table}[t]
\centering
\caption{Speedup scaling with prefix length}
\label{tab:speedup_scaling}
\small
\begin{tabular}{lccccc}
\toprule
Model & Prefix & Base (ms) & Ours (ms) & Speedup \\
\midrule
TinyLlama & 50 & 13.1 & 6.0 & 2.2$\times$ \\
TinyLlama & 200 & 14.8 & 4.0 & 3.7$\times$ \\
TinyLlama & 500 & 27.6 & 4.9 & 5.6$\times$ \\
TinyLlama & 1000 & 52.8 & 6.0 & 8.8$\times$ \\
TinyLlama & 1500 & 67.8 & 7.0 & \textbf{9.7$\times$} \\
\midrule
Mistral-7B & 500 & 234.2 & 35.0 & 6.7$\times$ \\
Mistral-7B & 1000 & 340.3 & 40.5 & 8.4$\times$ \\
Mistral-7B & 1500 & 450.8 & 46.6 & \textbf{9.7$\times$} \\
\bottomrule
\end{tabular}
\end{table}

Table~\ref{tab:speedup_scaling} shows absolute latency reduction.
With 1500-token prefix, latency drops from 67.8ms to 7.0ms on TinyLlama (\textbf{9.7$\times$} speedup).

\subsection{Real-World Dataset Validation}

To understand practical boundaries, we evaluate on ShareGPT-style conversations.

\begin{table}[t]
\centering
\caption{Real-world dataset evaluation}
\label{tab:sharegpt}
\small
\begin{tabular}{lccc}
\toprule
Experiment & Speedup & Token Reuse & Finding \\
\midrule
Multi-turn (short ctx) & 0.90$\times$ & 43.9\% & Limited benefit \\
Long-context QA & 1.02$\times$ & 66.1\% & Breakeven \\
\bottomrule
\end{tabular}
\end{table}

Table~\ref{tab:sharegpt} reveals an important finding: \textbf{\deltacache{} provides limited benefit for short-context multi-turn conversations}.
The cache management overhead exceeds computation savings when prefixes are short and new content per turn is substantial.
This motivates deployment primarily in long-prefix scenarios.

\subsection{Eviction Policy Comparison}

\begin{table}[t]
\centering
\caption{Eviction policy comparison (TinyLlama-1.1B)}
\label{tab:eviction}
\small
\begin{tabular}{lccc}
\toprule
Policy & Hit Rate & Speedup & Overhead \\
\midrule
LRU & 99.5\% & 2.12$\times$ & Low \\
LFU & 99.5\% & 2.20$\times$ & Low \\
Composite & 99.5\% & 2.15$\times$ & Medium \\
Tiered (Layer-aware) & 99.5\% & 2.11$\times$ & Medium \\
Adaptive (Attention) & 99.5\% & \textbf{2.20$\times$} & Medium \\
\bottomrule
\end{tabular}
\end{table}

Table~\ref{tab:eviction} compares eviction policies.
All achieve similar performance ($\sim$99.5\% hit rate) because high prefix sharing means eviction rarely triggers.
LFU and Adaptive show marginally better speedup.

\subsection{Memory Usage Analysis}

\begin{table}[t]
\centering
\caption{GPU memory overhead analysis}
\label{tab:memory}
\small
\begin{tabular}{lccccc}
\toprule
Model & Base & Ours & Overhead & Speedup & Eff. \\
\midrule
TinyLlama-1.1B & 2303 & 2808 & 505 MB & 2.40$\times$ & 4.9$\times$/GB \\
Mistral-7B & 14491 & 16893 & 2402 MB & 5.12$\times$ & 2.2$\times$/GB \\
\bottomrule
\end{tabular}
\end{table}

Table~\ref{tab:memory} shows GPU memory overhead.
\deltacache{} adds 17--22\% memory overhead for cache storage, achieving 2.2--4.9$\times$ speedup per GB overhead (memory efficiency).
The overhead is practical for production deployment.

\subsection{Long Context Scaling}

\begin{table}[t]
\centering
\caption{Speedup scaling with context length (TinyLlama-1.1B)}
\label{tab:longctx}
\small
\begin{tabular}{ccccc}
\toprule
Context & Base (ms) & Cached (ms) & Speedup & Scaling \\
\midrule
501 & 35.2 & 17.6 & 2.00$\times$ & -- \\
1001 & 54.3 & 16.3 & 3.32$\times$ & +1.32$\times$ \\
1501 & 69.8 & 16.7 & 4.18$\times$ & +0.86$\times$ \\
1801 & 90.0 & 17.4 & \textbf{5.17$\times$} & +0.99$\times$ \\
\bottomrule
\end{tabular}
\end{table}

Table~\ref{tab:longctx} demonstrates that speedup scales approximately \textbf{2.43$\times$ per 1K additional tokens}, confirming the theoretical quadratic relationship.
Cached query latency remains nearly constant ($\sim$17ms) while baseline latency grows linearly.

\section{Discussion}

\paragraph{Optimal Application Scenarios}
Based on our experiments, \deltacache{} is most effective for:
\begin{itemize}
    \item \textbf{Long prefix ($>$500 tokens)}: RAG with long documents, few-shot with many examples, detailed system prompts
    \item \textbf{High query volume}: Same document/context queried multiple times
    \item \textbf{Code completion}: Fixed large context with varying completions
\end{itemize}

\paragraph{When NOT to Use \deltacache{}}
\begin{itemize}
    \item \textbf{Short prefix ($<$100 tokens)}: Cache overhead exceeds savings
    \item \textbf{Growing context}: Multi-turn conversations where each turn adds substantial new content
    \item \textbf{Single-use prompts}: No prefix sharing opportunity
\end{itemize}

\paragraph{Comparison with Serving Systems}
Unlike vLLM and SGLang which operate at the serving layer, \deltacache{} is a model-agnostic library that can be integrated with any transformer.
This provides flexibility for custom deployments and research experimentation.

\paragraph{Limitations}
Current limitations include: (1) single-request focus without batched inference optimization, (2) CPU offloading adds latency for GPU retrieval, (3) no support for dynamic prefix modification.

\section{Conclusion}

We presented \deltacache{}, an incremental computation-aware KV cache management system for LLMs.
Our key results:
\begin{itemize}
    \item \textbf{5.13$\times$} average speedup on Mistral-7B (fp16) across targeted benchmarks
    \item \textbf{13.43$\times$} maximum speedup on 2K-token contexts in baseline comparisons
    \item Speedup scales at \textbf{2.43$\times$ per 1K tokens} with memory overhead of only 17--22\%
    \item \textbf{100\% correctness} at fp16 precision verified
\end{itemize}

Crucially, \textbf{scenario matching is the key factor}: \deltacache{} excels in fixed-prefix workloads (RAG, few-shot, code completion) achieving 3--13$\times$ speedup, while multi-turn conversations with growing context show limited gains.
Ablation studies confirm that speedup is strongly correlated with prefix length, with 500+ tokens being the practical threshold.

\deltacache{} provides a practical, model-agnostic solution for production LLM deployments with long shared prompt prefixes.
Future work includes batched inference optimization, learned eviction policies, and deeper integration with serving frameworks.

\bibliographystyle{icml2026}
\bibliography{references}

%%%%%%%%%%%%%%%%%%%%%%%%%%%%%%%%%%%%%%%%%%%%%%%%%%%%%%%%%%%%%%%%
\newpage
\appendix

\section{Implementation Details}
\label{app:impl}

\paragraph{Memory Management}
Default thresholds: GPU 80\%, CPU 70\%.
Proactive offloading begins at 70\% GPU utilization to maintain headroom.
Block size: 64 tokens for efficient memory alignment.

\paragraph{Layer Weight Parameters}
Sigmoid parameters: $k=5$ (steepness), $\tau=0.3$ (transition point).
This gives early layers (0--30\%) weight $\approx$0.18, late layers (70--100\%) weight $\approx$0.88.

\paragraph{Attention Tracking}
EMA decay $\gamma=0.9$ for attention score averaging.
Default importance weights: $w_a=0.5$ (attention), $w_f=0.3$ (frequency), $w_r=0.2$ (recency).

\paragraph{RoPE Position Handling}
For models with Rotary Position Embeddings, we explicitly compute absolute position IDs:
\texttt{position\_ids = torch.arange(prefix\_len, prefix\_len + delta\_len)}.
This ensures cached and freshly computed KV states are numerically equivalent.

\section{Additional Experimental Results}
\label{app:results}

\begin{table}[h]
\centering
\caption{Query count scaling (cache amortization)}
\small
\begin{tabular}{cccc}
\toprule
Queries & Overall & Cached & Amortization \\
\midrule
5 & 2.14$\times$ & 2.84$\times$ & Low \\
10 & 2.63$\times$ & 3.32$\times$ & Medium \\
20 & 2.71$\times$ & 2.82$\times$ & Good \\
50 & 2.98$\times$ & 2.80$\times$ & High \\
\bottomrule
\end{tabular}
\end{table}

\begin{table}[h]
\centering
\caption{Multi-turn conversation evaluation}
\small
\begin{tabular}{lcccc}
\toprule
Model & Turns & Latency (ms) & Hit Rate & Token Reuse \\
\midrule
TinyLlama & 2 & 15.23 & 50.0\% & 30.7\% \\
TinyLlama & 3 & 16.75 & 66.7\% & 51.6\% \\
TinyLlama & 4 & 16.40 & 75.0\% & 61.4\% \\
\midrule
Mistral-7B & 2 & 132.31 & 50.0\% & 31.2\% \\
Mistral-7B & 3 & 130.35 & 66.7\% & 50.6\% \\
Mistral-7B & 4 & 130.56 & 75.0\% & 61.5\% \\
\bottomrule
\end{tabular}
\end{table}

\section{Benchmark Scenario Details}
\label{app:scenarios}

\paragraph{Code Completion}
Prefix: Large Python file with class definitions, imports, and documentation ($\sim$1450 tokens).
Queries: Function completion requests at different points in the file.

\paragraph{RAG Knowledge Bases}
Three knowledge bases tested:
(1) Python standard library documentation,
(2) Machine learning tutorial content,
(3) REST API reference documentation.
Each prefix $\sim$650 tokens, queries are natural language questions.

\paragraph{Few-shot Sentiment}
5 labeled examples (positive/negative movie reviews) as prefix ($\sim$280 tokens).
Test inputs: Novel movie review snippets for classification.

\paragraph{Long System Prompt}
Detailed chatbot instruction prompt defining persona, capabilities, and response format ($\sim$480 tokens).
Queries: Diverse user questions.

\end{document}
